% Haz de 14 grados del radar sobre un balde de 30 cm (escala: 1 m = 4 cm)
% Superficie del agua en y = 1.2; borde del balde en y = 1.4; medio ancho del balde 0.6
\begin{tikzpicture}[
    balde/.style={draw=gray!70!black, very thick},
    sensor/.style={fill=gray!70!black, minimum width=0.6cm, minimum height=0.35cm, inner sep=0pt},
    cota/.style={<->, >=Stealth, gray!60!black},
    rotulo/.style={font=\footnotesize, align=center}
]
    % ---------- Caso 1: sensor a 1 m del agua ----------
    \begin{scope}
        \fill[agua!25] (-0.6,0) rectangle (0.6,1.2);
        \draw[balde] (-0.6,1.4) -- (-0.6,0) -- (0.6,0) -- (0.6,1.4);
        \fill[radar!25] (0,5.2) -- (-0.49,1.2) -- (0.49,1.2) -- cycle;
        \draw[radar, thick] (-0.49,1.2) -- (0,5.2) -- (0.49,1.2);
        \node[sensor] at (0,5.38) {};
        \draw[cota] (1.0,5.2) -- node[right, rotulo] {1 m} (1.0,1.2);
        \node[rotulo, below] at (0,-0.15) {Haz de 25 cm\\\textcolor{agua!80!black}{\textbf{entra en el balde}}};
    \end{scope}

    % ---------- Caso 2: sensor a 2 m del agua ----------
    \begin{scope}[xshift=5cm]
        \fill[agua!25] (-0.6,0) rectangle (0.6,1.2);
        \draw[balde] (-0.6,1.4) -- (-0.6,0) -- (0.6,0) -- (0.6,1.4);
        \fill[radar!25] (0,9.2) -- (-0.98,1.2) -- (0.98,1.2) -- cycle;
        \draw[radar, thick] (-0.98,1.2) -- (0,9.2) -- (0.98,1.2);
        \node[sensor] at (0,9.38) {};
        \fill[red] (-0.6,1.4) circle (2.2pt) (0.6,1.4) circle (2.2pt);
        \draw[cota] (1.5,9.2) -- node[right, rotulo] {2 m} (1.5,1.2);
        \node[rotulo, below] at (0,-0.15) {Haz de 49 cm\\\textcolor{red!80!black}{\textbf{toca los bordes}}};
    \end{scope}

    % Ángulo del haz
    \node[rotulo, radar] at (5,9.95) {14°};
\end{tikzpicture}
